\cotefile{157-5}
\scanpage{149}
\mspage{65}
% archive sheet number: 148

\unsure{dont|donc} il devient évident à présent qu'il faut le faire de toutes façons. Je vais quand même \unsure{énoncer|résumer} le positif de mes perplexités :

\uline{Proposition-per\deleted{\unsure{p}}plexité} 4.3.1. Soient $M,N$ deux catégories pseudo-topos, d'où $M\boxtimes N$ comme dans 4.1. Soient $S\subset M$, $T\subset N$ deux petites sous-catégories strictement génératrices, et considérons l'inclusion pleinement fidèle correspondante

\diagram{\[
\begin{tikzcd}[ampersand replacement=\&,column sep=large]
\underset{\substack{\shortparallel\,\mathrm{def}\\\underline{\unsure{\mathrm{Hom}|\mathrm{Hum}}}^{!!}(M^{\unsure{o|0}}\times\unsure{N|N^o},\mathrm{Ens})}}{M\boxtimes N} \arrow[r,hook,"\alpha"'] \& S^\wedge\boxtimes T^\wedge \approxeq (S\times T)^\wedge \arrow[l,dashed,bend right=15,"{\varphi?}"']
\end{tikzcd}.
\tag*{(4.3.1.1)}
\]}

\inserted{Considérons} les conditions \unsure{suivantes|suiv.} \deleted{sont équivalentes}

a) $M\boxtimes N$ est un pseudo-topos

b) Pour $x\in\mathrm{Ob}\,M$, $y\in\mathrm{Ob}\,N$, $x\boxtimes y\in\mathrm{Ob}\,M\boxtimes N$ \unsure{existe}, i.e. le foncteur (commutant aux petites $\underleftarrow{\lim}$)

\[
F\longrightarrow F(x,y),\qquad M\deleted{\unsure{\times}}\boxtimes N\longrightarrow\mathrm{Ens}
\tag*{(4.3.1.2)}
\]

est représentable.

b$'$) Comme b), avec $x\in\mathrm{Ob}\,S$, $y\in\mathrm{Ob}\,T$.

c) Le foncteur $\alpha$ \deleted{\unsure{admet un adj.}} \inserted{(lequel commute aux petites $\underleftarrow{\lim}$)} admet un adjoint à gauche $\varphi:(S\times T)^\wedge\longrightarrow M\boxtimes N$.

On a alors les implications

\deleted{Démonstration}

\[
b\Longrightarrow b'\Longleftrightarrow c\mathrel{\deleted{\unsure{<}}}\Longrightarrow a.
\tag*{(4.3.1.3)}
\]

De plus, si $c$ est satisfait, on a \deleted{\unsure{précis|préci}} : pour $s\in S^\wedge$, $t\in T^\wedge$, \deleted{$\unsure{\rho}(s,t)$ par sa \unsure{définition}} considérons $(s,t)$ \deleted{$\unsure{\in}\, S^\wedge\boxtimes T^\wedge\simeq(S\times T)^\wedge$}\inserted{$\in(S\times T)\subset(S\times T)^\wedge\approxeq S^\wedge\boxtimes T^\wedge$}, \deleted{qui} s'explicite \deleted{\unsure{au}} comme

\margin{\diagram{\begin{tikzcd}[ampersand replacement=\&,column sep=small,row sep=small]
S\times T\arrow[r,"\mathrm{pr}_2"]\arrow[d,"\mathrm{pr}_1"] \& T\\
S \&
\end{tikzcd}}}

\[
\left\{\begin{aligned}
(s\deleted{\unsure{\times}},t)&=\mathrm{pr}_1^{\unsure{*|\alpha}}(s)\times\mathrm{pr}_2^{*}(t)
\deleted{\unsure{=(s,t)}}\inserted{\simeq s\overset{S^\wedge,T^\wedge}{\boxtimes}t}\in\mathrm{Ob}\,(S\times T)^\wedge\\
\text{i.e.}\qquad(s\deleted{\unsure{\times}},t)(s',t')&=s(s')\times t(t')=\mathrm{Hom}_{S}(s',s)\times\mathrm{Hom}_{T}(t',\unsure{t})\\
&=\mathrm{Hom}_{S\unsure{\times}T}((s',t'),(s,t)).
\end{aligned}\right.
\tag*{(4.3.1.4)}
\]

\scanpage{150}
\mspage{66}
% archive sheet number: 149

on a

\[
s\overset{(M,N)}{\boxtimes}t\simeq\varphi(s\deleted{\unsure{\times}},t),
\tag*{(4.3.1.4)}
\]

\margin{\uline{NB} On a bien sûr $(s,t)=s\overset{S^\wedge,T^\wedge}{\boxtimes}t$}

et $\varphi(S\times T)$ \deleted{est une} \inserted{$=\{s\boxtimes t\mid s\in S,t\in T\}$} est une partie strictement génératrice de $M\boxtimes N$. \deleted{Je ne \unsure{connais}}

\uline{NB} Je ne connais pas d'exemple où \deleted{\unsure{toutes|les}} les conditions envisagées ne soient satisfaites, mais doute pourtant à présent que même la plus faible le soit toujours, i.e. que $M\boxtimes N$ soit un pseudo-topos dès que $M$ et $N$ le sont.

\uline{Démonstration} \uline{de} 4.3.1. L'implication $b\Rightarrow b'$ est tautologique. L'implication $b'\Longleftrightarrow c$ résulte du

\noindent\begin{tabular}{@{}|p{.93\linewidth}@{}}
\uline{Lemme} 4.3.2. Soient $\unsure{\mathcal E|E},\deleted{\unsure{S}}X$ des \inserted{$\unsure{U|\mathcal U}$-}catégories, avec $X$ petite, $\mathcal E$ stable par petites $\underrightarrow{\lim}$,

\[
\alpha:\mathcal E\longrightarrow X^\wedge
\]

un foncteur. Pour que $\alpha$ admette un adjoint à gauche $\varphi$, il faut et il suffit que pour tout $x\in\mathrm{Ob}\,X$, le foncteur

\[
\left\{\begin{aligned}
\unsure{\zeta|\xi}&\longmapsto \mathrm{Hom}_{X^\wedge}(x,\alpha(\zeta))\deleted{\unsure{=}}=\alpha(\zeta)(x)\\
\mathcal E^{\unsure{o|\circ}}&\longrightarrow\mathrm{Ens}
\end{aligned}\right.
\]

soit représentable.
\end{tabular}

C'est évidemment nécessaire. La suffisance résulte du \deleted{\unsure{le fa}} fait que $X\subset X^\wedge$ est stricte%
\scanpage{151}%
\mspage{67}%
% archive sheet number: 150
ment générateur dans $X^\wedge$ i.e. que \unsure{tt} objet $\deleted{\unsure{E}}F$ de $X^\wedge$ est limite inductive $F=\underset{I}{\underrightarrow{\lim}}\,F_i$, les $F_i$ \unsure{ds|dans} $X$ (on prend $I=X_{/F}$), d'où

\[
\begin{aligned}
\deleted{\mathrm{Hom}(F\unsure{,}}\ \mathrm{Hom}(F,\alpha(\unsure{\zeta|\xi}))
&\simeq\underset{I}{\underleftarrow{\lim}}\ \underbrace{\mathrm{Hom}(F_i,\alpha(\zeta))}_{\mathrm{Hom}(\unsure{\rho|\varphi}(F_i),\zeta)}\\
&\simeq\mathrm{Hom}(\underrightarrow{\lim}\,\rho(F_i),\zeta),
\end{aligned}
\]

donc le \unsure{foncteur} $\zeta\mathrel{\unsure{\mapsto|\vdash}}\mathrm{Hom}(F,\alpha(\zeta))$ est représenté par $\rho(F)\overset{\mathrm{def}}{=}\underrightarrow{\lim}\,\rho(F_i)$.

Pour prouver $c\Longrightarrow a$, \deleted{il suffit de prouver que c) implique} \inserted{\unsure{Pour} l'existence} des \inserted{petites} $\underrightarrow{\lim}$\inserted{ dans $M\boxtimes N$, et l'existence d'une partie \unsure{strictmt} génératrice}, on invoque le lemme général

\noindent\begin{tabular}{@{}|p{.93\linewidth}@{}}
\uline{Lemme} 4.3.3. Soit $\alpha:\unsure{\mathcal E|E}\hookrightarrow\unsure{\mathcal F|F}$ un foncteur \unsure{pl} fidèle, admettant un adjoint à gauche \inserted{$\rho$}. \inserted{\textcircled{a}} Si les $\underrightarrow{\lim}$ d'un type donné $I$ existent dans $\mathcal F$, elles existent dans $\mathcal E$, et se calculent comme
\end{tabular}

\[
\underset{I\deleted{\unsure{(\mathcal E)}}}{\underrightarrow{\lim}}_{\!(\mathcal E)}\,\zeta_i
\simeq\rho\bigl(\underset{I}{\underrightarrow{\lim}}_{\!\mathcal F}\,\alpha(\zeta_i)\bigr)
\tag*{(4.3.3.1)}
\]

\inserted{\textcircled{b} Si $U$ est une partie \inserted{(petite)} str. génératrice de $\mathcal F$, $\rho(U)$ est une petite p. str. gén. dans $\mathcal E$.}

En effet, \unsure{si} $\zeta$ désigne le deuxième \unsure{mb}, et $y$ un objet variable de $\unsure{\zeta|\mathcal E}$, on aura

\[
\begin{aligned}
\mathrm{Hom}_{\mathcal E}(\zeta,y)&\overset{\mathrm{adj}}{\simeq}
\mathrm{Hom}_{\mathcal F}\bigl(\underset{I}{\unsure{\underline{\lim}|\underrightarrow{\lim}}}_{\!\mathcal F}\,\alpha(\zeta_i),\alpha(y)\bigr)\\
&\simeq\underset{I}{\underleftarrow{\lim}}\ \underbrace{\mathrm{Hom}_{\mathcal F}(\alpha(\zeta_i),\alpha(y))}_{\simeq\mathrm{Hom}_{\mathcal E}(\zeta_i,y)\ \text{\unsure{car} }\alpha\ \text{pl. fid.}}
\end{aligned}
\]

\scanpage{152}
\mspage{68}
% archive sheet number: 151

\[
\simeq\underset I{\underleftarrow{\lim}}\ \mathrm{Hom}_{\unsure{\mathcal E|E}}(\unsure{\zeta|\xi}_i,y),
\]

ce qui prouve que $\zeta$ est une limite inductive des $\zeta_i$ dans $\mathcal E$.

% slanted
\margin{Pour \unsure{prouver} b), on note que la pl. fid. de $\alpha$ implique que $\rho$ est \unsure{essentielt|un foncteur} \unsure{surjectif} ; comme $\rho$ commute aux $\underrightarrow{\lim}$, et que tt objet de $\unsure{\mathcal F|F}$ est \deleted{\unsure{de}} \unsure{de} la forme $\underset I{\underrightarrow{\lim}}\,x_i$, les $x_i$ dans $U$, on \unsure{aura} $\rho(x)=\underset I{\unsure{\underline{\lim}|\underrightarrow{\lim}}}\,\rho(x_i)$, les $\rho(x_i)$ dans $\rho(U)$, \unsure{et} on gagne.}

\deleted{Comme} \inserted{Ceci prouve que (moyennant c))} $M\boxtimes N$ est stable par \inserted{petites} $\underrightarrow{\lim}$, donc il reste à prouver que $M\boxtimes N$ admet une petite famille strictement génératrice. Cela \unsure{résultera|résulte} du \deleted{Lemme \unsure{n°} \unsure{su-}} \inserted{b) du dit lemme,}

\inserted{puisque $S\times T$ engendre str\textsuperscript{t} $(S\times T)^\wedge$. Cela prouve en} \deleted{\unsure{tions} de la proposition, que je vais prouver à présent.} même temps \unsure{la dernière assertion|les dernières assertions} de la proposition, sous réserve de prouver \deleted{Je vais \unsure{maintenant}} la formule (4.3.1.4), \deleted{en désignant \unsure{par} \unsure{$u=s\boxtimes t$}.}

Il faut donc prouver un iso fonctoriel en $F\in\mathrm{Ob}\,M\boxtimes N$

\[
\mathrm{Hom}(\deleted{\unsure{F},}\rho(s,t),F)\simeq F(s,t),
\]

mais par adjonction on a

\[
\begin{aligned}
\mathrm{Hom}(\deleted{\unsure{F},}\rho(s,t),F)&\simeq\mathrm{Hom}(\deleted{\alpha(F),\unsure{(s,t)}}\inserted{(s,t),\alpha(F)})=\underset{\substack{\shortparallel\\F(s,t),}}{\alpha(F)(s,t)}
\end{aligned}
\]

\deleted{ok. Pour voir que $\rho(S\times T)$ est strictement génératrice dans $M\boxtimes N$, on note que $S\times T$ est strictement génératrice dans $(S\times T)^\wedge$, et que, $\alpha$ étant pl. fid., \deleted{son adjoint à g. $\rho$ est \unsure{tjs?|ess.t} surjectif, et}}

\deleted{Reste à prouver $a\Longrightarrow c$. Mais si $M\boxtimes N$ est un ps. topos, alors tt \inserted{\unsure{contr.}} foncteur sur $M\boxtimes N$ qui}

\scanpage{153}
\mspage{69}
% archive sheet number: 152

\uline{Proposition} 4.3.4. Soit $M$ une $\unsure{U|\mathcal U}$-catégorie. Les conditions suivantes sont équivalentes

a) $M$ est une catégorie $U$-pseudo-topos, i.e. $M$ est stable par $U$-limites inductives, et admet une \inserted{$U$-}petite ss-catégorie \inserted{$S$ str.} génératrice.

b) Il existe une \deleted{\unsure{famille|foncteur}} $U$-petite catégorie $S$, et un foncteur $\unsure{i}:M\hookrightarrow S^\wedge$ \uline{pleinement fidèle}, ayant un adjoint à gauche.

\uline{Dém} On a $a\Rightarrow b$, car il suffit de prendre pour $S$ \inserted{dans b)} celui spécifié dans a), l'hyp. sur $S$ implique que le foncteur canonique $M\unsure{\longrightarrow|-}\deleted{\unsure{M}}S^\wedge$ composé de $M\mathrel{\mathop{\hookrightarrow}\limits^{\deleted{\unsure{h}}}}M^\wedge\mathrel{\mathop{\unsure{-|\longrightarrow}}\limits^{\mathrm{restr.}}}S^\wedge$ est pleinement fidèle. Qu'il admette un adj. à g. résulte du lemme 4.3.2.

L'implication $b)\Rightarrow a)$ résulte du lemme 4.3.3.

\uline{Corollaire} \inserted{4.3.5.} Soit $M$ une $U$-catégorie. Pour que $M$ soit un $U$-pseudo-topos \deleted{\unsure{$\zeta|\xi$}} \uline{accessible} (i.e \unsure{tous|tout} \unsure{dont|dans} les objets sont accessibles), il f. et s. qu'il existe $S$ dans $\mathrm{Cat}$ et un foncteur pleinement fidèle et \uline{accessible} $M\mathrel{\mathop{\unsure{\lhook\joinrel-|\hookrightarrow}}\limits^\alpha}S^\wedge$, admettant un adjoint à gauche $\rho$.

\scanpage{154}
\mspage{70}
% archive sheet number: 153

\uline{Dém} \deleted{\unsure{Il}} suffit de \deleted{\unsure{noter|voir}} \inserted{noter} que si $S,\alpha$ est comme dans 4.3.4.\deleted{\unsure{a}}b), alors $\alpha$ est accessible \unsure{sss|ssi} \unsure{les|le} $\rho(S)$ \deleted{\unsure{$(s\in\mathrm{Ob}$}} est formé d'objets accessibles de $M$ ; comme cette partie est strict\textsuperscript{mt} génératrice, on sait que cela implique que \uline{tt} objet de $M$ est accessible. (Cf SGA 4 I 9.9.)

\deleted{\uline{\unsure{Corollaire}} \unsure{4.3.6}}

Je vais essayer de donner un énoncé plus complet :

\uline{Théorème} 4.3.6. Soit $M$ une $\unsure{U|\mathcal U}$-catégorie Conditions équivalentes

a) $M$ stable par petites $\underrightarrow{\lim}$, et elle admet une petite sous-catégorie pleine, \inserted{$S$} strictement génératrice \unsure{$S$} et formée d'objets accessibles de $M$

a$'$) $M$ est une \inserted{catégorie $U$-}pseudo topos \inserted{\deleted{\unsure{faiblement}}} accessible, i.e. satisfait a), avec en plus \uline{\deleted{\unsure{tout|tt}}} élément de $M$ (pas seulement ceux de $S$) accessible dans $M$

b) Il existe une petite catégorie $S$, et un foncteur $\alpha:M\unsure{\lhook\joinrel-|\hookrightarrow} S^\wedge$ \uline{pleinement fidèle} \uline{et accessible}, tel que $\alpha$ admette un adjoint à gauche.

\scanpage{155}
\mspage{71}
% archive sheet number: 154

\inserted{\deleted{\unsure{d}}c) $M$ est stable par petites $\underrightarrow{\lim}$, est accessible, \unsure{il|et} admet une filtration cardinale $(\mathrm{Filt}^{\pi}M)_{\pi\geq\pi_0}$}

\deleted{\unsure{c}}d) \deleted{\unsure{$M$}}Il existe un cardinal $\pi_{\unsure{0|1}}$, une petite catégorie $M_{\pi_0}$ stable par limites inductives de cardinal $\leq\pi_0$, \deleted{\unsure{un foncteur}} \inserted{telles qu'on} ait une équivalence de catégories

\[
M\overset{\approx}{\longleftarrow}\mathrm{Ind}_{\pi_0}\,M_{\pi_0},
\tag*{(4.3.6.1)}
\]

où le deuxième membre désigne la sous-catégorie pleine de $\mathrm{Ind}(M_{\pi_0})$ formée des \deleted{systèmes \unsure{$I$} \unsure{ind.}} ind-objets $(X_i)_{i\in I}$ de $M_{\pi_0}$ tels que $I$ soit grand devant $\pi_0$.

\uline{Démonstration de} 4.3.6. On a déjà $a'\Rightarrow a\Rightarrow b\Rightarrow a$, i.e l'équivalence \unsure{a},a$'$,b, soit $\mathrm{PTA}$ cette condition. Nous allons prouver les implications

% slanted
\margin{Pour une démonstration complète (celle donnée ici \unsure{coule|cale} à partir d'un certain moment) cf § 4.6 (\unsure{p|pp} 139--148).}

\[
\mathrm{PTA}\Longrightarrow c\Longrightarrow d\Longrightarrow\mathrm{PTA},
\]

ce qui établira le théorème.

\inserted{Prouvons} \uline{$\mathrm{PTA}\Longrightarrow c$.} Soit $S$ une petite sous-catégorie pleine \inserted{str.} génératrice. \deleted{Il est} \unsure{Ainsi|Alors}, \deleted{immédiat que si nous \unsure{a}}

\scanpage{156}
\mspage{72}
% archive sheet number: 155

pour tout $x$ dans $M$ on a

\[
x=\underset{s\in S_{/x}}{\underrightarrow{\lim}}\ s
\tag*{(4.3.6.2)}
\]

cf SGA 4 I \unsure{7|?}.2. \unsure{(i)|bis}). Soit $\pi_0=\mathrm{card}\,\mathrm{Fl}(S)$, \deleted{et soit} et pour tt cardinal $\pi\geq\pi_0$, soit

\[
\begin{aligned}
\mathrm{Filt}^{\pi}M^{\deleted{\unsure{\pi}}}
&=\text{sous-catégorie pleine de }M\text{ formée}\\
&\phantom{{}={}}\text{des objets }x\text{ tels que }\mathrm{card}(\mathrm{Fl}\,S_{/x})\leq\pi.
\end{aligned}
\tag*{(4.3.6.3)}
\]

Je veux prouver que \deleted{\unsure{ceci définit}} \inserted{l'on \unsure{trouve|s'trouve} ainsi} une filtration cardinale sur $M$, i.e. que les conditions a) b) c) de SGA 4 I 9.12. sont satisfaites.

C'est prouvé dans loc. cit. 9.13. sous des \deleted{conditions} hypothèses légèrement différentes, la condition d'accessibilité de $M$ y est remplacée par celle \deleted{que} (beaucoup moins sympa !) que les familles de flèches épimorphiques strictes dans $M$ sont épimorphiques strictes universelles. Cette condition n'est utilisée que dans la démonstration de \inserted{la} condition \unsure{b)}, à savoir la stabilité de $\mathrm{Filt}^\pi M$ par \deleted{\unsure{petites}} \inserted{$\underset{I}{\underrightarrow{\lim}}$} de cardinal $\leq\pi$ (sans avoir d'ailleurs à supposer $I$ filtrante, ni a fortiori grand devant $\pi_0$).

Il faut d'ailleurs dire que la définition de $\mathrm{Filt}^\pi M$ donnée dans loc cit. n'est pas tout à fait la $\unsure{=|\simeq}$ que (4.3.6.3),

\scanpage{157}
\mspage{73}
% archive sheet number: 156

mais c'est

\[
\begin{aligned}
\mathrm{Filt}^{\prime\,\pi}M
&=\text{ss-catégorie pleine de }M\text{ formée}\\
&\phantom{{}={}}\text{des objets }x\text{ tels qu'il existe}\\
&\phantom{{}={}}\text{une famille épimorphique}\\
&\phantom{{}={}}\text{stricte }\ x_i\xrightarrow{\alpha_i}x\ \text{de cardinal}\\
&\phantom{{}={}}\leq\pi,\ \text{avec }x_i\in S\quad\forall\,i\in I.
\end{aligned}
\tag*{(4.3.6.3)}
\]

\unsure{Aussitôt|Aisément|A priori} on a

\[
\mathrm{Filt}^\pi M\subset\mathrm{Filt}^{\prime\,\pi}M,
\]

mais dans loc. cit on prouve que pour $x$ dans $\mathrm{Filt}^{\prime\,\pi}M$, on a

\[
\mathrm{card}\,\mathrm{Fl}(S_{/x})\leq\pi^{\pi_0},
\]

d'où

\[
\mathrm{Filt}^\pi M\subset\mathrm{Filt}^{\prime\,\pi}M\subset\mathrm{Filt}^{\pi^{\pi_0}}M
\]

donc si $\pi=\pi^{\pi_0}$, p. ex. $\pi=2^c$ avec $c\geq\pi_0$, donc on trouve

\[
\mathrm{Filt}^\pi M=\mathrm{Filt}^{\prime\,\pi}M.
\]

Donc les deux filtrations ont l'air $\pm$ équivalentes. \deleted{Pour n'avoir pas à refaire du travail de SGA 4 I 9.13, je vais essayer de travailler avec la filtration définie dans loc. cit., plutôt qu'avec celle indiquée d'abord.}

Mais la démonstration donnée de la \deleted{condition b) de SGA 4 I 9.12.} stabilité de $\mathrm{Filt}^{\prime\,\pi}$ \deleted{fait} semble vraiment faire appel de façon essentielle à la condition d'universalité des familles épimorphiques strictes, (\deleted{\unsure{donc|dans}} \inserted{\unsure{opéra-|à prio-}} \unsure{attention|portion}, \unsure{\}} car les hypothèses faites \deleted{\unsure{déjà}} impliquent l'existence de petites $\underleftarrow{\lim}$).

Donc \unsure{il|zut}, me faut \unsure{ben} oui, \unsure{mal gré|malgré|mal que} refaire le travail !

\scanpage{158}
\mspage{74}
% archive sheet number: 157

Condition a) : $\mathrm{Filt}^\pi$ équivalente à une petite catégorie : c'est immédiat, puisque \unsure{par} (4.3.6.2) $\mathrm{Filt}^\pi$ est formé des $\underset{I}{\unsure{\underline{\lim}|\underrightarrow{\lim}}}$ \inserted{dans $M$} d'objets de $S$ (petite catégorie), avec $I$ telles que $\mathrm{card}\,\mathrm{Fl}\,I\leq\pi$.

Condition b) : $\mathrm{Filt}^\pi$\textsuperscript{*} stable par $\underset{I}{\underrightarrow{\lim}}$, $\mathrm{card}\,\mathrm{Fl}\,I\leq\pi$. Soit donc

% slanted
\margin{\unsure{prendre|prend} $I$ filtrant \unsure{gr.|et gr.} devant $\pi_0$ \unsure{je ne} le \unsure{prouve} que dans ce cas, hélas !}

\[
(x_i)_{i\in I},\ \deleted{\unsure{\text{ou}|\text{un}}}\qquad I\mathrel{\mathop{\longrightarrow}\limits^{\unsure{\varphi|\rho}}_{i\,\unsure{\mapsto|\vdash}\,x_i}}M
\]

un syst. inductif dans $M$, avec des $x_i$ dans $\mathrm{Filt}^\pi$. Considérons le produit fibré

\diagram{\[
\begin{tikzcd}[ampersand replacement=\&,column sep=large,row sep=large]
I\arrow[d,dash,"\varphi"'] \& S_{/I}=J\arrow[l,"\text{cofibré}"']\arrow[d]\arrow[dr,dashed,"\psi"] \& \\
M \& S_{/M}\arrow[l,"\text{but}"',"\text{Cat-cofibrée}"]\arrow[r,"\text{source}"] \& S\unsure{\lhook\joinrel-|\subset|\hookrightarrow}M
\end{tikzcd},
\tag*{(4.3.6.4)}
\]}

où $J$ est Cat-cofibré sur $I$, sa fibre en $i\in\mathrm{Ob}\,I$ étant isom. à $S_{/x_i}$, donc $\mathrm{Card}\,\mathrm{Fl}\,\unsure{I_i|J_i}\leq\pi$. Comme par hypothèse $\mathrm{card}\,\mathrm{Fl}\,I\leq\pi$, on en conclut que $\mathrm{card}\,\mathrm{Fl}\,J\leq\pi^3=\pi$,

[NB L'ens. des flèches \inserted{$v:j_0\longrightarrow j_1$} de $J$ au dessus d'une flèche $u:i_0\longrightarrow i_1$ de $I$ est ``contenu'' dans l'ens. produit $\mathrm{Ob}\,I_{i_0}\times\mathrm{Fl}\,I_{i_1}$, par l'application $v\mathrel{\unsure{\mapsto|\vdash}}(\text{source }v,v':u_*(\unsure{i_0|j_0})\longrightarrow j_1)$ donc de cardinal majoré par $\pi^2$.] \deleted{\unsure{Soit} \unsure{$x=\lim$}} On

\scanpage{159}
\mspage{75}
% archive sheet number: 158

a alors

\[
x\overset{\mathrm{def}}{=}\underset{I}{\underrightarrow{\lim}}\ \underset{\rule{.3pt}{2pt}\rule{1.2em}{.3pt}\rule{.3pt}{2pt}}{x_i}
=\underset{J}{\underrightarrow{\lim}}\ \underset{\substack{\rule{.3pt}{2pt}\rule{1.2em}{.3pt}\rule{.3pt}{2pt}\\\shortparallel\\\psi(j)}}{\unsure{\zeta|\xi}_{\unsure{j|s}}},
\]

\margin{\noindent\begin{tabular}{@{}|p{.92\linewidth}@{}}
On utilise la propriété de distributivité des $\underrightarrow{\lim}$, rel. à un foncteur Cat-cofibrant \diagram{\begin{tikzcd}[ampersand replacement=\&,row sep=small,column sep=small]
J\arrow[r]\arrow[dr,"\psi"'] \& I\\
\& M
\end{tikzcd}} explicitée dans
\end{tabular}}

où cette fois la limite est une limite d'objets de $S$, \inserted{prise dans $M$.} Il s'ensuit que $x\in\mathrm{Filt}^{\prime\,\pi}$, ce n'est pas tout à fait $x\in\mathrm{Filt}^{\pi}$, \deleted{Mais sauf \unsure{er}} [\unsure{ça|ce} implique $x\in\mathrm{Filt}^{\pi^{\pi_0}}$, \deleted{\unsure{et}} donc $x\in\mathrm{Filt}^{\pi}$ si $\pi=\pi^{\pi_0}$, p. ex. $\pi=2^c$ avec $c\geq\pi_0$]

Mais c'est le moment d'utiliser le fait que les objets de $S$ sont accessibles, et on prendra $\pi_0$ tel que les objets de $S$ soient $\pi_0$-accessibles, et cette fois on prendra quand-même $I$ \deleted{filtrante} et grande devant $\pi_0$ (pour nous borner à ce qui est exigé dans la condition b) de loc. cit.) On doit prouver que $\mathrm{Card}(\mathrm{Fl}\ S_{/x})\leq\pi$ (ce qui signifie $x\in\mathrm{Filt}^{\pi}$). Mais comme les objets de $S$ sont $\pi_0$-accessibles et $I$ grand devant $\unsure{k_0|\pi_0}$, on voit aussitôt que

\[
S_{/x}\simeq\underset{i}{\underrightarrow{\lim}}\ S_{/x_i}
\]

d'où

\[
\mathrm{Fl}\,S_{/x}\simeq\underset{i}{\underrightarrow{\lim}}\ \mathrm{Fl}\,S_{/x_i}
\]

donc

\[
\mathrm{Card}\,\mathrm{Fl}\,S_{/x}\leq
\underbrace{(\mathrm{Card}\ I)}_{\leq\pi}
\underbrace{\bigl(\underset{i}{\mathrm{Sup}}\ \mathrm{Card}\,\mathrm{Fl}\,S_{/x_i}\bigr)}_{\leq\pi}
\leq\pi^2=\pi
\]

\unsure{qed|cqfd}.

\scanpage{160}
\mspage{76}
% archive sheet number: 159

% slanted
\margin{\unsure{?} \uline{Démonstration \unsure{concluante|conclut}} \unsure{en vérité|en réalité} que si $\pi^{\pi_0}=\pi$ !!}

Prouvons c), i.e. que pour tt $\pi'>\pi\geq\pi_0$ et \deleted{tout} $x\in\operatorname{Ob}\operatorname{Filt}^{\pi'}$ on peut écrire

\[
\begin{aligned}
x=\varinjlim_I x_i\qquad & I\ \text{grand devant }\pi\\
&\text{les }x_i\text{ dans }\operatorname{Filt}^{\pi}\\
&\operatorname{Card} I\leq\pi'^{\pi}.
\end{aligned}\tag*{(4.3.6.5)}
\]

La démonstration se fait comme dans \deleted{Pour ceci, on peut supposer que} loc. cit. (p. \unsure{151|159}--152), en changeant légèrement la définition de $I$, comme désignant l'ens des sous-catégories pleines \inserted{$i$} de $S_{/x}$ qui sont filtrantes, de cardinal $\leq\pi$ \uline{et de plus grandes devant $\pi_0$}. (On pourra supposer \inserted{les} $S_{/x}$ filtrantes, en \deleted{\unsure{ajoutant}} \unsure{augmentant}\deleted{\unsure{t}} $S$ de façon à ce que $S$ soit stable \inserted{dans $M$} par \deleted{petites} $\varinjlim$ finies, donc les catégories $S_{/x}$ aussi, ce qui implique qu'elles sont filtrantes. Il est

% slanted
\margin{\unsure{?} Il n'est pas sûr que $I$ soit grand devant $\pi$, sauf si on suppose $\pi^{\pi_0}=\pi$ !}

On a donc prouvé $\mathrm{PTA}\Longrightarrow\mathrm{c})$,

\uline{Prouvons} $\mathrm c\Longrightarrow\mathrm d$. Soit \deleted{\unsure{Soit} $c\geq\pi_0$, et} \deleted{$c\geq\pi_0$, et soit}

\[
\pi_1=2^c\quad\text{avec }c\geq\pi_0.
\]

On sait alors que

\[
\begin{aligned}
\operatorname{Filt}^{(\pi_1)}M=M_{\pi_1}\overset{\mathrm{def}}{=}\ &\text{sous-catégorie \inserted{pleine de $M$} des}\\
&\text{objets }\pi_1\text{-accessibles}\\
&\text{de }M
\end{aligned}
\]

(loc. cit 9.16), ce qui implique notamment que $M_{\pi_1}$ est stable par $\varinjlim$ de cardinal $\leq\pi_1$, \inserted{en vertu de loc cit. 9.9. (p\deleted{\unsure{ag}} 145).}

\deleted{En effet, si $x=\varinjlim_{i\in I}x_i$, $\operatorname{card}\operatorname{Fl}I\leq\pi$, et $x_i\in M_{\pi_1}$, alors le foncteur \inserted{$h_x:M\unsure{\text{---}|\longrightarrow}\mathrm{Ens}$} représenté \unsure{par} \unsure{$h$|$x$} est la $\varprojlim_I$ des foncteurs $h_{x_i}$ ; comme ils sont $\pi_1$-accessibles, il en est de $=$ de $h_x$ (par loc. cit. 9.\unsure{6} (ii\unsure{i})).}

De plus, on

\scanpage{161}
\mspage{77}
% archive sheet number: 160

par loc. cit. \deleted{\unsure{§}}9.18. b)\inserted{ (p. 156)}, on a

\[
\operatorname{Ind}_{\pi_1}(M_{\pi_1})\xrightarrow{\approx}M,
\]

ce qui prouve d).

\uline{Prouvons $\mathrm{d})\Longrightarrow\mathrm{PTA}$} Par loc. cit \deleted{\unsure{§}}9.\unsure{18|8} a)\inserted{ (p. 156)}, \deleted{\unsure{on sait}} on voit que les objets de $M_{\pi_1}$ sont $\pi_1$-accessibles dans $\operatorname{Ind}_{\pi_1}(M_{\pi_1})$, \deleted{\unsure{donc}} d'autre part il est clair que $M_{\pi_1}$ est une sous-catégorie pleine strictement génératrice de \deleted{\unsure{$\mathcal C$|$M$}} $\operatorname{Ind}_{\pi_1}(M_{\pi_1})$, donc il reste à voir que $\operatorname{Ind}_{\pi_1}(M_{\unsure{1|\pi_1}})$ est stable par petites $\varinjlim$, i.e. le

\uline{Lemme} 4.3.\deleted{\unsure{6}}7. Soit $\unsure{M|A}_\pi$ une \inserted{\unsure{$U$|$\mathcal U$}-}catégorie stable par limites inductives de cardinal $\leq\pi$, où $\pi$ est un cardinal \inserted{infini} donné. Alors $\operatorname{Ind}_{\pi}M_{\pi}$ est stable par \uline{petites} limites inductives. De plus, \inserted{$M_\pi\longrightarrow M$} $M=\operatorname{Ind}_{\pi}\unsure{C|M_\pi}$ est solution d'un problème 2-universel, à \unsure{savoir} que \deleted{pour} \deleted{\unsure{$M$}} pour tte \deleted{\unsure{$U$-}}catégorie $N$ stable par petites limites inductives, le foncteur restriction \deleted{\unsure{ind}} induit une équivalence

\[
\underline{\operatorname{Hom}}_{!}(M,N)\xrightarrow{\approx}
\underline{\operatorname{Hom}}_{!\pi}(M_{\pi},N),
\tag*{(4.\unsure{3|8}.7.1)}
\]

où \unsure{le|les} \deleted{\unsure{prem}} \unsure{signe|signes} $!$ (resp. $!\pi$) désigne la sous-catégorie pleine de la catégorie $\underline{\operatorname{Hom}}$, formée des foncteurs qui commutent aux petites $\varinjlim$ (resp. aux limites inductives de cardinal $\leq\pi$).

\scanpage{162}
\mspage{78}
% archive sheet number: 161

% slanted
\margin{Mais \uline{\unsure{Attention|Att.}} $M_\pi\longrightarrow\operatorname{Ind}(M_\pi)$ ne \unsure{commute} \unsure{$\neq$|pas} \deleted{\unsure{pas|$\Rightarrow$}} à ces $\varinjlim$ !}

Prouvons d'abord que les $\varinjlim$ de cardinal $\leq\pi$ \uline{dans $M_\pi$} sont aussi des $\varinjlim$ \uline{dans $M$}\inserted{$=\operatorname{Ind}_\pi(M_\pi)$}. On a donc \deleted{dans $M_\pi$} $x$

\[
x\xleftarrow{\sim}\varinjlim_I x_i\qquad \text{\uline{\unsure{ds|dans}} } M_{\pi},
\]

i.e.

\[
\operatorname{Hom}(x,y)\xrightarrow{\sim}\varprojlim_I\operatorname{Hom}(x_i,y)\qquad
\text{si } y\in\operatorname{Ob} M_\pi\tag*{$(\ast)$}
\]

prouvons que cela reste vrai pour $y$ dans $M$ i.e. $y$ de la forme

\[
\begin{aligned}
y=\varinjlim_J y_j\qquad &J\ \text{grand devant }\pi\\
&\text{les } y_j\text{ dans } M_\pi.
\end{aligned}
\]

On a en effet

\[
\begin{aligned}
\operatorname{Hom}(x,y)&\simeq\varinjlim_J\operatorname{Hom}(x,y_j)
\overset{(\ast)}{\simeq}\varinjlim_J\varprojlim_I\operatorname{Hom}(x_i,y_j)\\
&\underset{\substack{\mathrm{SGA}\,4\,\mathrm I\\9.8}}{\simeq}\varprojlim_I\varinjlim_J\operatorname{Hom}(x_i,y_j)
\simeq\varprojlim_I\operatorname{Hom}(x_i,y),
\end{aligned}
\]

au total

\[
\operatorname{Hom}(x,y)\simeq\varprojlim_I\operatorname{Hom}(x_{\unsure{i|s}},y)
\]

qed. Pour prouver que les \inserted{petites} $\varinjlim$ \deleted{de type \uline{fini}} sont représentables dans $M=\operatorname{Ind}_\pi(M_\pi)$, comme elles le sont dans $\operatorname{Ind}(M_\pi)$ (SGA 4 I 8.9.5 \unsure{b|b)}\unsure{)}), en utilisant l'hypothèse que les $\varinjlim$ \uline{finies} sont, représentables dans $M_\pi$), il reste à voir que $\operatorname{Ind}_\pi(M_\pi)$ est stable dans $M_\pi$ pour ces limites.

\deleted{Soit donc $(X_j)_{j\in J}$ un système inductif dans $M$\inserted{$=\operatorname{Ind}_\pi(M_\pi)$}, $x$ sa limite dans $\operatorname{Ind}(M_\pi)$, identifié à un objet de $M_\pi^\wedge$, via $M_\pi\hookrightarrow\operatorname{Ind}_{\pi}(M_\pi)\hookrightarrow\operatorname{Ind}(M_\pi)\hookrightarrow M_\pi^\wedge$,}

\scanpage{163}
\mspage{79}
% archive sheet number: 162

4.4. \uline{Étude \unsure{de|des}} $\operatorname{Ind}_\pi(C)$ (pour $C$ \unsure{petite} catégorie stable par $\varinjlim$ de cardinal $\leq\pi$).

4.4.1. Finalement, la démonstration du lemme 4.3.7 n'est pas tellement évidente que je l'\unsure{anti}cipais, et il me faut lui consacrer des préliminaires techniques, que je \unsure{réunirai|vais réunir} ici. Dans tte la suite, $C$ représente une \inserted{\unsure{$U$|$\mathcal U$}-}catégorie équivalente à une petite catégorie.

\uline{Proposition 4.4.2.} Soit $F\in C^\wedge$ un foncteur ind-représentable, i.e tel que $C_{/F}$ soit filtrante. Soit d'autre part $i\longmapsto x_i$ un foncteur $\unsure{\varphi|q}\,I\longrightarrow C$, où $I$ est une petite catégorie filtrante, et \deleted{considérons} soit

\[
\alpha:x=\varinjlim_I{}_{\!(C^\wedge)}x_i\longrightarrow \deleted{\unsure{E|F}}F
\tag*{(4.4.2.1)}
\]

% slanted
\margin{La donnée de $\alpha$ équivaut aussi à celle d'une flèche de ind-objets de $(x_i)_{i\in I}\longrightarrow F$, $F$ étant considéré comme un Ind-objet \deleted{\unsure{constant}} de $C$ \unsure{indexé par} $C/F$.}

un système projectif de flèches $x_i\unsure{-|\longrightarrow} F$ (NB, la limite inductive $x$ est prise dans $C^\wedge$, et est calculée argument par \uline{argument}), \inserted{ce qui équivaut à} \deleted{de sorte que $\unsure{\varphi|q}\,I\longrightarrow C$ \unsure{se}} une factorisation $\varphi:I\longrightarrow C_{/F}$ de $I\longrightarrow C$ en

\diagram{\begin{tikzcd}[ampersand replacement=\&]
I \arrow[r,"\varphi_\alpha"] \& C_{/F} \arrow[r,hook,"\mathrm{can}.", "\unsure{\psi|+|r}"'] \& C.
\end{tikzcd}}

Ceci dit, pour que $\alpha$ (4.4.\unsure{2|?}.1) soit un isomorphisme, dans $\operatorname{Ind}'(C)\subset C^\wedge$ (ss-catégorie pleine des foncteurs ind-représentables sur $C$), il f. et s. que $\varphi_\alpha$ soit \uline{cofinal} (ou ce qui revient au $=$, $I$ étant filtrant) qu'il satisfasse les deux conditions habituelles

\[
\begin{aligned}
\mathrm{F1})\quad &\forall\ \unsure{\zeta|\xi}\in\operatorname{Ob} C_{/F},\ \exists i\in\operatorname{Ob} I\text{ et une flèche }\zeta\longrightarrow\varphi_\alpha(i)\\
\mathrm{F2})\quad &\forall\ \deleted{\unsure{z}}\zeta\in\operatorname{Ob} C_{/F},\ \text{et }i\in\operatorname{Ob} I,\text{ et tte double flèche}
\end{aligned}
\]

\scanpage{164}
\mspage{80}
% archive sheet number: 163

\diagram{\begin{tikzcd}[ampersand replacement=\&]
\unsure{\zeta|\xi} \arrow[r,shift left=.55ex] \arrow[r,shift right=.55ex] \& \varphi_\alpha(i),
\end{tikzcd}}

existe une flèche $i\overset{u}{\unsure{-|\longrightarrow}}j$ dans $I$ telle que $\varphi_\alpha(u)$ égalise la double flèche.

\uline{Dém} Si $\varphi_\alpha$ est cofinal, on en \unsure{conclut|déduit} par définition que $\varinjlim\deleted{\unsure{\psi|\varphi}}\varphi_\alpha\longrightarrow\varinjlim\psi\varphi_\alpha$ est iso, or cette flèche n'est autre que $\alpha$. Inversement, supposons que cette flèche soit un iso, donc $C_{/F}\xleftarrow{\sim}C_{/\varinjlim\deleted{\unsure{I}} x_i}\unsure{)}$, et il reste à voir que le foncteur can.

\[
I\longrightarrow C_{/\varinjlim_{(C^\wedge)}x_i}
\]

est cofinal, i.e. satisfait F1) et F2). La vérification est $\pm$ \unsure{tautologique}, en utilisant le fait que pour tt \deleted{\unsure{$\zeta$}}$y$ \deleted{\unsure{$\in$}} \inserted{$\leftarrow$} dans $C$ \deleted{(et \unsure{notamment un quelconque} \unsure{des termes} $x_j$, $j\in\operatorname{Ob} I$),}

\[
\text{on a}\qquad \operatorname{Hom}(y,\varinjlim_{C^\wedge}x_i)\xleftarrow{\sim}\varinjlim_I{}_{\mathrm{Ens}}\operatorname{Hom}(y,x_i),
\]

et du ``calcul'' des limites inductives filtrantes dans Ens. \unsure{$\mathsection$|$\square$}

\uline{Corollaire} 4.4.3. \inserted{Soit $\pi$ un cardinal \unsure{infini}.} Pour que $F$ soit dans $\operatorname{Ind}^{\unsure{+|'}}_{\pi}(C)$\newline
\begin{tabular}{@{}|p{.91\linewidth}@{}}
(image essentielle de $\operatorname{Ind}_{\pi}(C)$ dans $C^\wedge$) il f. et s. qu'il existe une catégorie ordonnée gr. de devant $\pi$ $I$, et un foncteur cofinal $I\longrightarrow C_{/F}$.
\end{tabular}

Il \deleted{faut donc} nous faut donc donner une condition, sur une catégorie $\unsure{J|\mathcal J}$ $(=C_{/F})$, pour qu'il existe une catégorie ordonnée $I$ \deleted{\unsure{grande}} devant $\pi$, et un foncteur cofinal $I\longrightarrow J$.

